\documentclass[12pt,a4paper]{article}
\usepackage{amsmath,amsthm,amsfonts,amssymb,amscd}
\usepackage{times}              
% Use Times New Roman
\usepackage{graphicx}           
% Enhanced support for images
\usepackage{float}              
% Improved interface for floating objects
\usepackage{booktabs}           
% Publication quality tables
\usepackage{xcolor}             
% Driver-independent color extensions
\usepackage{geometry}           
% Customize document dimensions
\usepackage{fullpage}           
% all 4 margins to be either 1 inch or 1.5 cm
\usepackage{comment}            
% Commenting
\usepackage{minted}             
% Highlighted source code. Syntax highlighting
\usepackage{listings}           
% Typeset programs (programming code) within LaTeX
\usepackage{lastpage}           
% Reference last page for Page N of M type footers.
\usepackage{fancyhdr}           
% Control of page headers and footers
\usepackage{hyperref}           
% Cross-referencing 
\usepackage[small,bf]{caption}  
% Captions
\usepackage{multicol}
\usepackage{tikz}               
% Creating graphic elements
\usepackage{circuitikz}         
% Creating circuits
\usepackage{verbatim}          
% Print exactly what you type in
\usepackage{cite}               
% Citation
\usepackage[us]{datetime} 
% Various time format
\usepackage{blindtext}
% Generate blind text
\usepackage[utf8]{inputenc}
\usepackage{array}
\usepackage{makecell}
\usepackage{tabularx}
\usepackage{titlesec}

\input {defs.tex}



\begin{document}

\textcolor{UM_Brown}{
\begin{minipage}{0.1\textwidth}
    \begin{flushleft}
        \includegraphics[height=3.5cm]{UM_Logo_VERT_CMYK.eps}
    \end{flushleft}
\end{minipage}
\begin{minipage}{0.8\textwidth}
    \begin{center}
        \textbf{\Large HVDC Transmission I}\\
        \vspace{5pt}
        Assignment 1 \\
        \vspace{20pt}
        \textit{Ramin Parvari} \\
        \vspace{5pt}
        \today
    \end{center}
\end{minipage}
\vspace{10pt}
\hrule
}



%%%%%%%%%%%%%%% NEW SECTION %%%%%%%%%%%%%%% 
\section*{Problem 1: Mathematical expressions}
The Fourier transform of the signal $x(t)$ is defined as
\begin{equation}
    X(f)=\int_{-\infty}^{\infty} x(t) e^{-j2\pi ft} \, \tu{d}t
    \label{eq: F trans}
\end{equation}
where $X(f)$ is the spectrum of the signal at frequency $f$. The time domain expression can then be reproduced by synthesizing all frequency content through inverse Fourier transform as follows.
\begin{equation}
    x(t)=\int_{-\infty}^{\infty} X(f) e^{j2\pi ft} \, \tu{d}f
    \label{eq: inv F trans}
\end{equation}

By taking the derivative of the equation (\ref{eq: inv F trans}) with respect to the time $t$, one can prove the derivative property of the Fourier transform as
\begin{equation}
    \frac{\tu{d}x}{\tu{d}t} \qquad \to \qquad (j2\pi f) X(f)
    \label{eq: deriv.}
\end{equation}



%%%%%%%%%%%%%%% NEW SECTION %%%%%%%%%%%%%%% 
\section*{Problem 2: How to include table?}

\begin{table}[H]
    \caption{Characteristics of the buck converter}
    \label{tab: Data}
    \centering
    \begin{tabular}{c|c|c}
    \toprule
     \bf{Input Voltage (V)}& \bf{Output Voltage (V)} & \textbf{Calculated Duty Cycle} \\
     \midrule
     40 & 20 & 0.5\\
     50 & 20 & 0.5 \\
     80 & 20 & 0.25 \\
     \bottomrule
    \end{tabular}
\end{table}



%%%%%%%%%%%%%%% NEW SECTION %%%%%%%%%%%%%%% 
\section*{Problem 3: How to include graphics?}
Images, figures, and photos are usually included in the context using the command \texttt{includegraphics}  inside \texttt{figure} environment.

\begin{verbatim}
    \begin{figure}[H]
        \centering
        \includegraphics[options]{filename}
        \caption{caption for the figure}
        \label{fig: figlabel}
    \end{figure}
\end{verbatim}

Figure \ref{fig: ElemntsPE} shows some elements that are frequently used in power electronic converters. This figure is adopted from \cite{FPE}.

\begin{figure}[H]
    \centering
    \includegraphics[width=15cm]{figures/elements.png}
    \caption{Elements of power electronic converters}
    \label{fig: ElemntsPE}
\end{figure}

For more information on how to use this package, please visit the following website: \\
\url{https://latexref.xyz/_005cincludegraphics.html}



%%%%%%%%%%%%%%% NEW SECTION %%%%%%%%%%%%%%% 
\section{Problem 4: How to cite?}
Paper \cite{Gole97} provides guidelines for modeling power electronics in electric power engineering application. If you do not include any reference please delete the reference section at the end of the document (see three lines of code before \texttt{\textbackslash end\{document\}})









%%%%%%%%%%%%%%% REFERENCES %%%%%%%%%%%%%%% 
\newpage
\bibliographystyle{IEEEtran}
\bibliography{References}








\end{document}
